\documentclass{article}

% Language setting
% Replace `english' with e.g. `spanish' to change the document language
\usepackage[spanish, es-nodecimaldot]{babel}
\usepackage[most]{tcolorbox}

% Set page size and margins
% Replace `letterpaper' with `a4paper' for UK/EU standard size
\usepackage[letterpaper,top=2cm,bottom=2cm,left=3cm,right=3cm,marginparwidth=1.75cm]{geometry}

% Useful packages
\usepackage{amsmath}
\usepackage{graphicx}
%\usepackage{subfigure}
\usepackage{caption}
\usepackage{subcaption}
\usepackage[colorlinks=true, allcolors=blue]{hyperref}

\newtcolorbox[auto counter, number within=section]{exercise}[1][]{
    colback=blue!5!white,
    colframe=blue!75!black,
    fonttitle=\bfseries,
    title=Ejercicio \thetcbcounter,
    #1 % Allows you to add custom options locally (like labels)
}




\title{Reinforcement learning II}
\author{Arturo Gil y Juan José Cabrera}

\begin{document}
\maketitle

\section{Introducción}
Con esta práctica se complementan los conceptos vistos en el área de Q-learning. En concreto, debemos comenzar recordando la principal limitación del algoritmo Q-learning: ``el problema de la alta dimensionalidad de los datos''. En la práctica anterior vimos claramente que el problema de aprendizaje se complicaba cuando las dimensiones de la tabla $Q(s, a)$ crecían. Por otra parte, también debemos notar que en la misma práctica de \textit{reinforcement learning} se resolvían problemas donde el estado y las acciones eran de naturaleza discreta.

\textbf{Resumen de la práctica}: en esta práctica se presentará inicialmente una solución de Q-learning  para el caso discreto, que utiliza una discretización directa del estado. El estudiante experimentará esta solución utilizando un simulador de  \textit{módulo lunar}. A continuación, se observará las limitaciones de Q-learning para resolver el problema (debido al problema de la su alta dimensionalidad). Se introduce en esta práctica una variación de esta solución que intenta evitar la estimación directa de una tabla Q: Deep Q Learning. En esta solución, se reemplaza la tabla Q por una red neuronal que se estimará en base al estado, acciones y recompensas en cada instante recibidas por el agente. En resumen, veremos estas dos técnicas y las aplicaremos al mismo problema:
\begin{itemize}
    \item Problema: entorno \textit{Lunar Lander}.
    \item Técnica 1: Q-learning discreto.
    \item Técnica 2: Deep Q-learning.
\end{itemize}


%\begin{verbatim}
%    https://medium.com/@coldstart_coder/dqn-algorithm-training-an-ai-to-land-on-the-moon-1a1307748ed9
  

 %   https://www.kaggle.com/code/mauriciofigueiredo/dqn-on-lunar-lander-rl-with-videoplay

  %  https://arxiv.org/pdf/1711.07478

    %https://web.stanford.edu/class/psych209/Readings/MnihEtAlHassibis15NatureControlDeepRL.pdf

    %file:///home/arturo/Baixades/Deep-Q-learning-explained-final.pdf
%\end{verbatim}

\section{Descripción del entorno}
El agente debe aprender a aterrizar un módulo lunar en la zona marcada con dos banderas blancas. El módulo es atraído hacia abajo por la gravedad. El estado inicial es aleatorio, pudiéndose encontrar el \textit{lander} con diferentes posiciones, velocidades y orientaciones iniciales. El objetivo consiste en aterrizar de forma segura en la zona indicada con una velocidad cercana a cero en el mínimo tiempo posible, mediante el accionamiento de los motores y con el mínimo consumo de combustible.

\subsection{Estado}
El estado está representado por las siguientes variables:
\begin{itemize}
    \item[1.] Posición horizontal $x$ $(m)$.
  \item[2.] Posición vertical $y$ $(m)$.
  \item[3.] Velocidad horizontal $v_x=\dot{x}$ $(m/s)$.
  \item[4.] Velocidad vertical $v_y=\dot{y}$ $(m/s)$.
  \item[5.] Ángulo $\theta$ $(rad)$.
  \item[6.] Velocidad angular $\omega=\dot{\theta}$ $(rad/s)$.
  \item[7.] Contacto de la pata izquierda $l_1$ (sí/no, 1/0).
  \item[8.] Contacto de la pata derecha $l_2$ (sí/no, 1/0).
\end{itemize}

Así pues, el estado en el instante $t$ se representa por $s_t=(x, y, v_x, v_y, \theta, \omega, l_1, l_2)$.

\subsection{Acciones}
Las acciones del agente son:
\begin{itemize}
    \item[ 1.] No hacer nada.
    \item[2.] Activa el motor principal.
  \item[3.] Activa el motor izquierdo.
  \item[4.] Activa el motor derecho.
\end{itemize}

 
\subsection{Estado inicial}
El \textit{lander} comienza en la parte central superior de la pantalla con una fuerza aleatoria aplicada a su centro de masas. Así pues, comienza con una aceleración aleatoria en cada episodio (pero con velocidad nula).

\subsection{Finalización del episodio}
El episodio finaliza si:
\begin{itemize}
    \item El \textit{lander} se estrella (el cuerpo del lander entra en contacto con la luna).
    \item El \textit{lander} sale de la pantalla (su coordenada $|x|>1$).
    \item El \textit{lander} no se mueve (tiene velocidad nula).
\end{itemize}


\begin{exercise}
    Intente aterrizar el módulo lunar. Ejecute el script \texttt{lunar\_lander\_test\_human.py} (práctica de  Q-learning II, parte A). Anote sus estadísticas sobre 20 intentos de aterrizaje. ¿Cuál es la recompensa media obtenida? ¿Se ha conseguido aterrizar en alguna ocasión?
\end{exercise}

\subsection{Recompensas}
En cada paso de simulación se recibe una recompensa. La forma en la que se recibe la recompensa es un poco diferente a lo que hemos visto hasta ahora. Para calcular la recompensa, gymnasium comienza planteando una valoración global $K_t$, basada en el vector de estado:
$$\text{K}_t = -100 \sqrt{x^2 + y^2} - 100 \sqrt{v_x^2 + v_y^2} - 100 \vert{}\theta\vert{} + 10 \cdot \text{l}_1 + 10 \cdot \text{l}_2$$

La posición $(x, y)$ se mide de forma relativa al centro de la plataforma de aterrizaje y el ángulo $\theta$ se mide respecto de la vertical. Los valores discretos $l_1$ y $l_2$ toman valores de 1 o 0 en función de si la pata toca el suelo.  La recompensa $r_t$ no se calcula de forma absoluta, sino como la diferencia en el valor de $K$ entre dos instantes de tiempo:
\begin{equation*}
\text{r}_t = \text{K}_t - \text{K}_{t-1} 
\end{equation*}

Sobre la recompensa $r_t$ se resta un valor $0.3$ si se usa el motor principal y de $0.03$ si se usan los impulsores laterales. Con esto, típicamente:
\begin{itemize}
    \item La recompensa aumenta si el módulo lunar se acerca al centro de la plataforma de aterrizaje (disminuye, en caso contrario).
    \item La recompensa aumenta cuando el lunar \textit{lander} se mueve más despacio (disminuye, si se mueve más rápido).
    \item La recompensa disminuye cuando el \textit{lander} se inclina más (aumenta cuando el \textit{lander} se acerca a $\theta = 0$ $ (rad)$).
    \item Se incrementa con 10 puntos con cada pata que entre en contacto con el suelo.
\end{itemize}

Dada la naturaleza acumulativa de la recompensa, tenemos que la recompensa total es: $\sum_{t=1}^Tr_t=\sum_{t=1}^T \text{K}_t - \text{K}_{t-1} = \text{K}_T - \text{K}_0$. Así pues, moverse desde la parte superior de la pantalla (con $\text{K}_0 \approx -100 \text{ a } -140$) hasta la plataforma de aterrizaje (con $\text{K}_T \approx +20$) da una recompensa total de +100 a +140 puntos por cada episodio finalizado correctamente. El episodio recibe una recompensa adicional de -100 o +100 puntos en caso de que se estrelle contra el suelo o aterrice de forma segura, respectivamente. Se considera que un episodio es una solución satisfactoria si consigue, por lo menos, 200 puntos.



Esta forma de diseñar la recompensa tiene, como principal ventaja, que se previene que el agente abuse de la obtención de recompensas. Debido que las recompensas depende de una diferencia calculada sobre una función global $K$, el agente no puede ganar infinitos puntos intentado volar en círculos cerca de la plataforma de aterrizaje: igual que se ganan puntos de recompensa acercándose al punto de aterrizaje, se pierden cuando el \textit{lander} se aleja.

\subsection{Recompensas (bis)}
Nos damos cuenta de que las recompensas en este entorno están diseñadas de una forma muy diferente al problema del \textit{Taxi} (o \textit{Cliff Walking}). Recordemos que en esos casos, básicamente, se daba una recompensa alta cuando el agente cumplía la misión y se daba una recompensa baja y negativa en cada paso del episodio. En este caso, la variación en la recompensa ya le da pistas al agente de qué es más efectivo para poder aterrizar el aparato en el lugar designado.

En el caso del lunar \textit{lander}, es importante fijarse en que el estado está formado por una serie de variables continuas. En este caso, nos damos cuenta de que, si diéramos una recompensa de, por ejemplo, $+100$ cuando se aterriza correctamente y una recompensa de $-1$ en cada paso del episodio, las probabilidades de que tuviéramos una secuencia de comandos sobre los motores que produjera un aterrizaje exitoso sería prácticamente nula. Y también al contrario, si el \textit{lander} se estrella después de 200 pasos del episodio, ¿cómo podemos determinar si se tomó una acción en un determinado momento que produjo que se descontrolara la aeronave? (por ejemplo, dar la vuelta a la aeronave hacia abajo, al haber accionado demasiado el motor derecho).


\subsection{Gravedad}
Extrañamente, encontraréis que el parámetro de gravedad que se indica por defecto en el entorno es $g=-10m/s^2$. Sobre la superficie de la luna, un objeto debería sufrir una aceleración de $g=-1,62m/s^2$. Este parámetro no es demasiado relevante, pues con cualquiera de estos valores se se trata de un entorno difícil de controlar. Si seleccionamos $g=-10m/s^2$ (valor por defecto en \textit{gymnasium}, realmente, estaremos simulando algo parecido al LLRV\footnote{\url{https://en.wikipedia.org/wiki/Lunar_Landing_Research_Vehicle}} (figura \ref{fig:LLRV}), utilizado por los intrépidos pilotos del programa \textit{Apollo} para ensayar las maniobras de aterrizaje sobre la luna (sí: ensayaban y aprendían en un entorno diferente... con diferente aceleración, claro está). También se permite introducir un parámetro de velocidad de viento en el entorno, que dificulta todavía más la tarea. 

\begin{figure}
    \centering
    \includegraphics[width=0.5\linewidth]{figures/LLRV.jpg}
    \caption{El Lunar Lander Research Vehicle}
    \label{fig:LLRV}
\end{figure}

\section{¿Un problema de control?}
En efecto, podríamos planificar una trayectoria desde cualquier estado del módulo lunar que lo llevara hasta la zona de aterrizaje. Se plantearía así el diseño de una función de control que calculará las acciones necesarias para seguir la trayectoria de referencia. La mayor parte de la teoría de control se basa en:
\begin{itemize}
    \item La estimación correcta de la planta $G(s)$ a controlar (el módulo lunar, en este caso).
    \item El cálculo de un controlador $R(s)$.
\end{itemize}
En esta práctica se persigue una solución que no implique la estimación de la planta $G(s)$ ni el cálculo explícito de un regulador $R(s)$.

\section{Q Learning por discretización}
En este apartado se presenta una solución para el aprendizaje del \textit{lunar} \textit{lander} mediante la discretización del estado y una estimación de la tabla $Q$ para esos valores discretizados. Cuando el entorno y el agente se describen mediante un estado $s_t$ y unas acciones $a_t$ discretas, es fácil almacenar los valores de $Q(s,a)$ en una tabla, como vimos en la práctica anterior. Recordemos, además, que en el caso del \textit{Taxi} y del \textit{Cliff Walking}, el número total de estados era reducido. En el caso del módulo lunar, el estado es continuo y las acciones discretas. Se propone discretizar el estado de la siguiente manera:
\begin{small}
\begin{verbatim}
    X_BINS = np.array([-0.4, -0.15, -0.05, 0.05, 0.15, 0.4])
    Y_BINS = np.array([0.05, 0.2, 0.5, 0.9])
    VX_BINS = np.array([-0.5, -0.2, -0.05, 0.05, 0.2, 0.5])
    VY_BINS = np.array([-0.8, -0.3, -0.1, 0.0, 0.2])
    ANGLE_BINS = np.array([-0.4, -0.1, -0.02, 0.02, 0.1, 0.4])
    ANG_VEL_BINS = np.array([-0.6, -0.2, 0.2, 0.6])
       
    def discretize_state(self, state):
        """Maps the 8D continuous state into a discrete tuple of bin indices."""
        x, y, vx, vy, angle, ang_vel, leg1, leg2 = state
        return (
            int(np.digitize(x, X_BINS)),
            int(np.digitize(y, Y_BINS)),
            int(np.digitize(vx, VX_BINS)),
            int(np.digitize(vy, VY_BINS)),
            int(np.digitize(angle, ANGLE_BINS)),
            int(np.digitize(ang_vel, ANG_VEL_BINS)),
            int(leg1),
            int(leg2),
        )
\end{verbatim}
\end{small}

La función de discretización asigna un número entero en función del \textit{bin} en el que caiga el valor de cada variable del estado $s_t=(x, y, v_x, v_y, \theta, \omega, l_1, l_2)$. Nótese que las variables que indican de contacto de las patas del módulo ya son discretas. 
Por ejemplo, considere que en un determinado instante el  estado $s_t$ tiene estos valores:
$s_t = (0.003, 1.41,  0.36, 0.32, -0.004, -0.081, 0, 1)$ $(m, m, m/s, m/s, rad, rad/s, [1], [1])$.
El estado se discretiza y representará por: 
$s^*_{t} = (3, 4, 5, 5, 3, 2, 0, 1)$.

%Por ejemplo, considere que $x=0.55m$ en un determinado instante, la función de discretización le asignará un valor de 7.


\begin{exercise}
    ¿Después de la discretización propuesta, cuántos estados diferentes se pueden codificar? Compare el resultado con el número de estados en el caso de los entornos Taxi y Cliff Walking.
\end{exercise}

En efecto, estamos discretizando de forma variable cada una de las partes del estado $s_t$. Por ejemplo, $x$ se discretiza en 7 \textit{bins} diferentes, $y$ en 5 intervalos diferentes, $v_x$ en 7, ... etc. Estas divisiones (\textit{bins}) se han elegido con mucha cautela, buscando reducir el número total de estados en la tabla Q. Con todo esto, en total tendremos un número total de posibles estados discretos considerablemente grande:
$$\text{Número de estados totales} = 7 \times 5 \times 7 \times 6 \times 7 \times 5 \times 2 \times 2=\mathbf{205.800}$$

\begin{exercise}
    Vaya al método \texttt{train} de la clase QLearningD y observe en qué puntos se hace la discretización y justifique por qué.
\end{exercise}

\subsection{Diccionarios}
Python permite guardar la tabla $Q$ en un diccionario en vez de en un array. Es importante pensar que se debe almacenar el valor de $Q(s_t, a_t)$: dado un estado $s_t$ hay que almacenar 4 posibles acciones.
En el caso de utilizar un array de numpy, es necesario hay que almacenar $\mathbf{205.800}\times4$ acciones totales. Si cada acción se guarda en un float de 64 bits (8 bytes), la cantidad de memoria total que se necesita es: $\mathbf{205.800}\times4\times8=6.585.600 \text{ bytes} \approx \mathbf{6,28 \text{ MB}}$.

En el código que se proporciona al estudiante, se usa la siguiente línea para definir la tabla $Q$:
\begin{verbatim}
    self.q_table = defaultdict(lambda: np.zeros(self.action_size))
\end{verbatim}
En esencia, en las líneas anteriores se define el uso de un diccionario de python. El diccionario de tipo {\textttdefaultdict} es especialmente rápido. Esto quiere decir que se usa un diccionario de python para almacenar el estado. La clave de cada elemento será una combinación del estado (un hash) y para cada clave se guardarán los 4 valores de la tabla, para cada acción. La sintaxis con \texttt{lambda: np.zeros(self.action\_size)} dice que se devuelva un array de 4 ceros cuando se intenta consultar un estado que no existe. El ahorro de memoria al usar defaultdict es importante, pues no todos los estados se visitan durante el aprendizaje o el test. Piense que seguramente no se visiten nunca algunas combinaciones extremas del estado. 

\section{Actividades}
Se proporciona al estudiante el código de la versión discretizada de Q-Learning. A continuación se proponen una serie de ejercicios.

\begin{exercise}
    Ejecute un entrenamiento de, al menos 5000 episodios (\texttt{lunar\_lander\_train.py} (Q-learning II, parte A). Utilice $\alpha=0.5$ y $\gamma=1.0$. Observará que, al acabar el entrenamiento, aparecen unos resultados parecidos a los de la figura \ref{fig:lunar_lander_first_train}. En la figura se observa:
    \begin{itemize}
        %\item La recompensa total de cada episodio.
        %item La recompensa media de 100 episodios (media móvil).
        \item La recompensa media del test inline cada 50 episodios.
        \item El valor de $\epsilon$ (100$\times$).
        %\item El número de estados visitados (1000$\times$).
    \end{itemize}
    ¿Qué conclusiones se pueden extraer?\\
    \textbf{IMPORTANTE}: observe que, tras el entrenamiento, se guarda un fichero pickle que contiene el diccionario de la tabla Q. Este fichero se leerá posteriormente en el script de test.
\end{exercise}
Un posible resultado del anterior script lo vemos en la figura \ref{fig:lunar_lander_first_train}.
Vemos que el aprendizaje es lento (vemos que la curva de la recompensa media sube poco a poco). También vemos que el aprendizaje no es del todo consistente: se obtiene, por ejemplo, una recompensa media de -150 en el episodio 800 de entrenamiento que baja a -250 en el episodio 1000.

%Vemos que la recompensa media sube correctamente en los primeros episodios, pero se queda atascada después. Nos fijamos, también, que los estados visitados son escasos, respecto del total de estados. Se propone hacer que $\epsilon$ decaiga más despacio.

\begin{figure}
    \centering
    \includegraphics[width=0.8\linewidth]{figures/lunar_lander_first_train.png}
    \caption{Resultado de un entrenamiento inicial del lunar lander (Q-learning discreto).}
    \label{fig:lunar_lander_first_train}
\end{figure}

\begin{exercise}
    Ejecute el método de test (script \texttt{lunar\_lander\_test.py}). Observe el comportamiento del agente. ¿Se consigue completar la tarea siempre?
\end{exercise}
Tras ejecutar el test en modo visual, observamos que en algunos casos se resuelve correctamente el episodio (la recompensa es mayor de 200) y en muchos casos no.

\subsection{Estrategia $\epsilon$-greedy}
En este caso se utiliza una estrategia para $\epsilon$ geométrica 25/75. Se pueden hacer variaciones sobre la ratio anterior, aunque el resultado sobre el aprendizaje es muy complejo de evaluar.
%En este caso, se plantea hacer una simplificación para la variación de $\epsilon$. Se plantea que durante el 80\% de los episodios, se haga una disminución lineal de $\epsilon$, para mantenerlo a un valor mínimo durante el resto de episodios.



\subsection{Manejo y visualización de resultados}
En el script de entrenamiento se ha incluido una función \texttt{save\_results} que guarda los resultados a un fichero \texttt{.json} (directorio \texttt{results}). Cada vez que se ejecuta el script de \texttt{train}, se guarda un nuevo fichero en función de los hiperparámetros $\gamma$ y $\alpha$. En cualquier momento podrá ejecutar el script \texttt{plot\_results} para observar los resultados de diferentes experimentos en un navegador web (abra el fichero .html). Esto nos va a permitir probar con diferentes valores de los hiperparámetros y observar qué diferencias hay en el aprendizaje.

\begin{exercise}
Haga pruebas con diferentes valores de $\alpha$ y $\gamma$ y $\delta$ (decay rate de $\epsilon$). Indique con qué parámetros se produce un mejor aprendizaje. Pruebe con valores de $\alpha$ entre 0.1 y 1.0. Varíe los valores de $\gamma$ entre 0.1 y 1.0. Utilice un número de episodios que no sea excesivo para poder realizar las pruebas (2000). \textbf{Una vez tenga una configuración que considere óptima, aumente el número de episodios e intente obtener un resultado final del aprendizaje con Q-learning tabular discretizado.}
\end{exercise}

Al realizar los experimentos anteriores, nos tenemos que dar cuenta de la gran variabilidad de los resultados: nos debemos fijar en la media de la recompensa obtenida, pero también en la varianza de los resultados. Así pues, es necesario:
\begin{itemize}
    \item Fijar unos parámetros de $\alpha$ y $\gamma$.
    \item Repetir varias veces el entrenamiento y test con esos parámetros.
    \item Guardar los datos para revisarlos posteriormente.
\end{itemize}

Este proceso es muy costoso computacionalmente. En la figura \ref{fig:results_discrete_total_episodes_15000_5_repetitions} se presenta el resultado de variar  $\alpha$ y $\gamma$ (con 5 repeticiones) con entrenamientos de 15.000 episodios (se empleó el script \texttt{lunar\_lander\_train\_opt.py}). Vemos que los parámetros $\alpha=0.5$ y $\gamma=1.0$ son los que permiten un aprendizaje mejor.  Con esto, para poder obtener un aprendizaje relevante, debemos entrenar, como mínimo, unos 15.000 episodios con estos parámetros. 

\begin{figure}
    \centering
    \includegraphics[width=0.8\linewidth]{figures/results_discrete_total_episodes_15000_5_repetitions.png}
    \caption{Recompensa media al variar $\alpha$ y $\gamma$ en entrenamientos de 15000 episodios (5 repeticiones).}
    \label{fig:results_discrete_total_episodes_15000_5_repetitions}
\end{figure}

\begin{exercise}
Una vez tenga una configuración de hiperparámetros que considere óptima, aumente el número de episodios de entrenamiento e intente obtener un resultado final del aprendizaje con Q-learning tabular discretizado. Pruebe el resultado con el script de test. Haga \texttt{render\_mode=None} y lo siguiente para poder probar exhaustivamente el conocimiento adquirido:
\begin{verbatim}
    res=qlearningd.test(total_episodes=500)
    plt.figure()
    plt.plot(res)
    plt.xlabel('Episodios')
    plt.ylabel('Recompensa')
    plt.title('Recompensa en episodios de test')
    plt.show()
\end{verbatim}
\end{exercise}
Si hemos hecho todo correctamente, se puede observar que algunos episodios se resuelven correctamente, con recompensas totales altas y cercanas a 200. Sin embargo, hay muchos casos en los que no se consigue aterrizar.

\subsection{Conclusiones}
El estudiante debe darse cuenta de que usar una técnica de Q-learning tabular no es la mejor solución para este problema. Dado que el estado es continuo, el hecho de haberlo discretizado ha hecho muy complejo el problema, requiriendo un gran número de iteraciones y episodios de entrenamiento para aumentar la inteligencia del agente. Nos hemos dado cuenta, también, de que es vital realizar una buena selección de los parámetros $\alpha$ y $\gamma$ para optimizar el aprendizaje del agente. Igualmente, es vital seleccionar correctamente los intervalos en los que se discretiza el estado. Nos damos cuenta de que también, al haber discretizado el problema, estamos perdiendo precisión a la hora de controlar el \textit{lander}.

\textbf{En resumen, y como conclusión principal, debemos ver que el problema del \textit{lunar lander} es muy complejo, debido a la gran cantidad de estados discretos que se han codificado y que el método de Q-learning discretizado está al límite de sus capacidades.} Después del entrenamiento, y durante las pruebas que realicemos, se puede ver que en algunas ocasiones el agente es capaz de aterrizar en la zona designada, pero esto no se consigue de forma consistente. Para mejorar los resultados, se debería hacer una discretización más fina del estado, que conllevaría la necesidad de ejecutar muchos más episodios de entrenamiento. En el apartado siguiente se presentará un método que permite conseguir resultados mucho mejores y con un menor esfuerzo computacional.


\section{Deep Q-learning}
En este apartado se presenta una técnica conocida como Deep Q-learning (DQN). En este método se evita la estimación de una tabla (o un diccionario) $Q(s,a)$ para estados discretos. En su lugar, se usa una red neuronal para estimar los valores de $Q(s,a)$ a partir del estado continuo.
El origen de esta técnica lo encontramos en el artículo titulado ``Playing Atari with Deep Reinforcement Learning'' \cite{MnihKSGAWR13}. En el paper, los autores demuestran que se puede construir un sistema de aprendizaje profundo capaz de aprender a utilizar 7 juegos de la Atari 2600 con resultados similares o mejores de los que obtendría un humano experimentado. El aprendizaje utiliza como entrada imágenes RGB sin procesar y obtiene como salida las acciones a realizar sobre el juego (que son aquellas que se podían realizar con un \textit{joystick}). La recompensa se lee directamente de la pantalla del videojuego. La idea principal de este método se presenta en la figura \ref{fig:dqn_agent}. El agente DQN está implementado mediante una red neuronal profunda (en concreto una red de tipo convolucional, CNN). Usando este paper como inspiración, se puede plantear una técnica similar que se aplicará al caso del entorno \textit{Lunar Lander}, donde el estado está representado mediante una serie de variables continuas (y no mediante una imagen de píxeles RGB). Las ideas principales de este método se listan, a continuación:
\begin{itemize}
    \item Se sustituye la tabla $Q$ por una red neuronal.
    \item La red neuronal se entrena para aproximar los valores $Q(s_t, a)$. Igual que antes, estos valores intentan aproximar la recompensa total que conseguirá el agente al realizar la acción $a$ cuando se encuentra en el estado $s_t$.
    \item La red neuronal se entrenará con una serie de muestras de entrada $(s_t, a_t)$ y salidas $Q(s_t, a)$. Las salidas que deberá aproximar la red son los mismos valores que se guardaban en una tabla $Q$:
    $$Q(s, a) = r_t+\gamma\max_a{Q(s_{t+1}, a)}$$
    Así pues, en el caso del \textit{Lunar Lander}, básicamente se trata de entrenar una red de regresión. Esta red tendrá de entrada el valor de $s_t$ y producirá como salida 4 valores de $Q(s_t, a)$. La figura \ref{fig:q_learning_mlp} muestra esta idea.
    \item Nótese que la anterior expresión es una versión de la Ecuación de Bellman con $\alpha=1$. El factor de aprendizaje no se introduce en la ecuación de Bellman, sino directamente como hiperparámetro de la red neuronal.
    \item Se busca que la red proporcione una generalización: debe ser capaz de estimar valores de  salida $Q(s_t, a_t)$, para valores de estado que no haya visto nunca antes. Esto es una ventaja enorme, si lo comparamos con el método que acabamos de ver, pues elimina el problema del aumento de complejidad con la discretización más fina del estado.
\end{itemize}

La red de regresión deberá aproximar 4 valores de $y$ a partir de la entrada $X=(s_t)$ (estado). Los valores de \textit{target} a estimar son los valores de la ecuación de Bellman: 
$$y=Q(s_t, a)=r_t+\gamma\max_a Q(s_{t+1}, a)$$

\textbf{Importante}: la red neuronal $Q$ aproxima una salida $y=Q(s_t, a)$ diferente para cada una de las 4 posibles acciones a tomar. 

Para mejorar la convergencia de la red neuronal a sus pesos óptimos, se utiliza un \textit{buffer} de datos (denominado \textit{mini-batch}). Así, en este \textit{batch} se almacenan la experiencia anterior del agente en términos de $(s_t, a_t)$. De este almacén temporal se escogen de forma aleatoria un número de muestras en cada iteración para el cálculo de los pesos de la red neuronal.

%El error en la estimación lo definimos como:
%$$e(s_t, a_t, r_t, s_{t+1}, w_t) = r_t
%+\gamma \max_a{Q(s_{t+1}, a_{t+1}, w_t)}$$

%Donde $w_t$ hace referencia a los pesos de la red neuronal. Se tendrá un error cuadrático medio:
%$$MSE = E[e^2]$$

%Finalmente, la red neuronal deberá actualizar sus pesos %$w$ de esta manera:
%$$w \leftarrow w +\alpha\cdot e \nabla Q(s_t, a_t, w_t)$$

\begin{figure}
    \centering    
    \includegraphics[width=0.8\linewidth]{figures/dqn_agent.png}
    \caption{Un esquema sencillo de la técnica Deep Q Learning.}
    \label{fig:dqn_agent}
\end{figure}

\begin{figure}
    \centering    
    \includegraphics[width=0.8\linewidth]{figures/q_learning_mlp.png}
    \caption{La tabla $Q$ se aproxima por una red MLP con 4 salidas (la dimensión de las acciones) y 8 entradas (la dimensión del estado).}
    \label{fig:q_learning_mlp}
\end{figure}


\subsection{Implementación}
En relación con la implementación que se proporciona al estudiantado, cabe decir:
\begin{itemize}
    \item La red neuronal se implementa usando las librerías de \texttt{scikit-learn}, pues se han estudiado y trabajado durante prácticas anteriores. En este caso, se usa una red \texttt{MLPRegressor} con 8 entradas y 4 salidas.
    \item Se pueden encontrar fácilmente en internet otras implementaciones basadas en las librerías \texttt{torch}, que permiten hacer uso de una GPU para el entrenamiento. No son estrictamente necesarias para resolver el problema en este caso.
    \item Se proporcionan otras clases que resultan de ayuda para poder implementar de forma sencilla y amable el algoritmo: principalmente un \textit{buffer} temporal de estados y acciones pasadas y una clase \texttt{EpsilonGreedyGeom} que se usó en una práctica anterior.
\end{itemize}


\subsection{Pasos del algoritmo}
Resumimos el algoritmo, a continuación, en una serie de pasos. 
\subsubsection{Inicialización}
Comenzamos creando dos redes neuronales con pesos iniciales $w$ aleatorios.
    \begin{itemize}
        \item $Q(s_t, a)$: que llamaremos \textbf{red neuronal online}.
        \item $\hat{Q}(s_t, a)$: que llamaremos \textbf{red neuronal target}.  
    \end{itemize}  
Las dos redes neuronales se usarán para aproximar vectores de estado-acción $(s_t, a_t)$ a valores de $Q$. Se fija inicialmente $\epsilon=1.0$ (estrategia $\epsilon$-greedy). Este valor se irá decrementando según transcurran los episodios, pasando el agente desde comportamiento aleatorio a realizar acciones meditadas según la red neuronal Q (que aproxima la tabla Q). Se crea un buffer para almacenar vectores $(s_t, a_t, r_t)$ que almacenan la experiencia realizada en episodios pasados y que se usará para estimar la red neuronal.

Se usan dos redes neuronales para estabilizar el proceso de estimación de los pesos de la red online. Este problema se origina por la siguiente circunstancia: dada la entrada $X=(s_t)$, se desea ajustar los pesos de la red para que produzcan 4 salidas con los valores $y=Q(s_t, a) = r_t + \gamma\max_a Q(s_{t+1}, a)$ (un valor Q para cada posible acción), de tal manera que podamos elegir la acción más correcta como la correspondiente al valor máximo de los 4 $Q(s_t, a)$. \textbf{Es importante darse cuenta de que $y$ depende, a su vez, de una estimación anterior de la red para calcular $Q(s_{t+1}, a)$ y este carácter cíclico no es adecuado para la estimación de los pesos de la red neuronal.} 

\subsubsection{Cálculo de acciones}
Durante el entrenamiento se usa una estrategia $\epsilon$-greedy para la selección de una acción (igual que en el método anterior). Así, con una probabilidad $\epsilon$ se elige una acción aleatoria $a_t=rand(S)$ (sobre todas las posibles acciones en el entorno). Por otra parte,  se realiza una acción meditada (explotación) con probabilidad $1-\epsilon$. La acción en este caso será la que maximice los valores de $Q$: 

$$a_t = \arg \max_a{Q(s_t, a)}$$.





\subsubsection{Aplicamos la acción y guardamos la experiencia}
Se ejecuta la acción $a_t$ en el entorno y se recibe una recompensa $r_t$ y el agente pasa al siguiente estado $s_{t+1}$. Se guardan estos datos en un buffer $(s_t, a_t, r_t, s_{t+1})$ de donde se pueden recuperar más tarde.

\subsubsection{Muestreo}
Se muestrean de forma aleatoria datos del buffer (denominado \textit{mini-batch}). Estos conjuntos de datos se usan para entrenar la red Q (online).

\subsubsection{Cálculo de los targets $y$}
Para cada vector $(s_t, a_t, r_t, s_{t+1})$ en el mini-batch se calcula un valor $y$, que es el que deberá aproximar la red de regresión. El valor de $y$ será:
\begin{equation*}
y = 
\begin{cases}
    r_t & \text{si el episodio termina en ese paso}, \\
    r_t + \gamma \max_a{\hat{Q}(s_{t+1}, a)}   & \text{en cualquier otro caso } 
\end{cases}
\end{equation*}

La razón de hacerlo de esta manera es para asegurar que los estados terminales tienen una recompensa alta durante el \textit{training}. Este valor alto de recompensa podría propagarse a estados anteriores, favoreciendo así el aprendizaje.

\subsubsection{Cálculo de la función de pérdida}
Se calcula una función de pérdida estándar entre los valores $y$ que es necesario aproximar y los resultados predichos por la red online $Q(s, a)$. Esto es:

$$ L=(Q(s,a) - y)^2 $$

Nótese que los valores de $y$ están calculados con la red $\hat{Q}$. Esta red es lenta y se actualiza solamente cada $N$ iteraciones.

\subsubsection{Actualización de la red online}
Se actualizan los parámetros de la red online $Q(s,a)$ usando los datos de error del minibatch utilizando SGD para minimizar la función de pérdida. 


\subsubsection{Actualización de la red $\hat{Q}$}
Cada $N$ pasos, se copian los pesos de la red $Q$ a la red $\hat{Q}$.

\subsubsection{Iteración}
Se itera el proceso hasta llegar a una convergencia. Esta convergencia se detecta por tener, de forma consistente, recompensas mayores de 200. Así pues, se repite el proceso hasta que podemos asegurar que se está obteniendo un aprendizaje correcto del proceso. En nuestro caso, no hacemos ninguna comprobación y observaremos el resultado del aprendizaje a través de una serie de gráficas.

\subsection{¿Por qué se usan dos redes neuronales?}
Hemos visto que se usan dos redes diferentes: $Q$ (online) y $\hat{Q}$ (red neuronal target). El propósito de la red $\hat{Q}$ es calcular los targets $y = r_t + \gamma \max_{a} \hat{Q}(s_{t+1}, a)$. Es un hecho conocido que se producen inestabilidades en la estimación si intentamos que la red $Q$ converja, mientras usamos $Q$ para estimar los valores a los que queremos que converja (es un problema de tipo ``pescadilla que se muerde la cola'').

\textbf{En resumen, la red neuronal target $\hat{Q}$ simplemente se mantiene constante durante unos episodios para estabilizar la estimación de $Q$ a través de la función de pérdida.}

\subsection{Pseudo-código}
Se da, a continuación, un pequeño trozo en pseudo-código para resumir todo el algoritmo anterior. También, el estudiante puede leer la clase \texttt{QLearningDQN} que implementa este método de aprendizaje por refuerzo.

\begin{small}
    \begin{verbatim}
        def train():
            Crear Q(s,a) (q_net, online network)
            Crear Q^(s, a), target_net
            Crear B (replay_buffer)
            for episode in range(total_episodes):
               at = greedy_epsilon()
               s', reward = self.env.step(s, action)
               # almacenar en buffer
               B.push([s, a, s', r])
               train_networks(episode)
               
        def greedy_epsilon():
            p = rand
            if p < epsilon:
                a=env.rand_action
            else:
                a = argmax Q(s,a) (a=argmax q_net.predict(s))
            return a
            
        def train_networks():
            # se actualiza la red online Q_net y
            # la red target
            se calculan los valores de y como
            y = rt (caso terminal)
            y = r_t + gamma max_a{Q^(s_{t+1}, a)}

            se actualiza Q(s,a) con el gradiente
            e=Q(s,a)-y
            cada N iteraciones guardamos Q
            Q^ = Q
    \end{verbatim}    
\end{small}

\section{Actividades}
Se proporciona una implementación funcional de esta técnica en la clase \texttt{QLearningDQN}. En las siguientes actividades el/la estudiante podrá observar y entender las principales líneas de código de la implementación. También se modificarán los parámetros principales del algoritmo para entender el funcionamiento del mismo. Nos interesará observar el progreso del aprendizaje del agente durante el entrenamiento. \textbf{Principalmente nos interesará ver qué ocurre al variar el parámetro $\gamma$ (\textit{discount factor}) y el tamaño de la red neuronal. También centraremos el interés en el tamaño del buffer y la frecuencia con la que se actualiza la red target (se hace $\hat{Q}=Q$).}

El resultado del algoritmo lo podemos evaluar en dos fases:
\begin{itemize}
    \item Durante el entrenamiento ejecutaremos tests inline para evaluar el aprendizaje. 
    \item Durante la fase de test, ejecutaremos episodios y calcularemos la recompensa total. Nos interesa, además, calcular la recompensa media y la varianza de esta recompensa: es interesante que el agente obtenga buenos resultados (y consiga aterrizar), pero, además, queremos que no se obtengan resultados desastrosos (y el \textit{lander} se estrelle).
\end{itemize}

Así pues, en resumen, estamos interesados en modificar los hiperparámetros del algoritmo y ver qué resultados se obtienen.

\subsection{Train/test}

\begin{exercise}
    Abra el script \texttt{lunar\_lander\_train.py} (práctica de Q-learning 2, parte B). Observe qué datos se guardan al finalizar el entrenamiento.
    Abra el script \texttt{lunar\_lander\_test.py} (práctica de Q-learning 2, parte B). Observe qué datos se leen en el método de test.
\end{exercise}

Nos damos cuenta de que, en este caso, al finalizar el entrenamiento no se guarda la tabla $Q$, sino que se almacenan los pesos de la red neuronal online. Estos pesos se guardan usando la librería de python \texttt{pickle} y tienen una extensión \texttt{.pkl}. Recordamos que esta red neuronal aproxima la tabla $Q$ (de hecho, hemos estado llamando $Q$ también a la red neuronal). Por otra parte, en el método de entrenamiento, es posible cargar los pesos de la red para continuar el entrenamiento desde donde se dejó. Igual que en la técnica anterior, se pueden observar los resultados del aprendizaje al finalizar el entrenamiento en una figura de \texttt{matplotlib}. 


\begin{exercise}
    Abra el script \texttt{lunar\_lander\_train.py} (práctica de Q-learning 2, parte B). Ejecute un entrenamiento de 1000 episodios con $\gamma=0.99$ y $N_{target}=0.25$. Quiere que se calculen tests \textit{inline} cada 100 episodios. Desea que cada test incluya 50 episodios. Observe los resultados de la gráfica.
\end{exercise}

Un resultado posible se muestra en la figura  \ref{fig:train_dqn_example}. Podemos observar que el valor de $\epsilon$ va disminuyendo durante los episodios. También observamos que los tests inline tienen resultados mejores. Cada punto de la gráfica corresponde con la media de varios tests que se hacen durante el entrenamiento.

\begin{figure}
    \centering
    \includegraphics[width=0.8\linewidth]{figures/train_discrete_q_learning_300000_episodes.png}
    \caption{Resultado de un entrenamiento de 30.000 episodios.}
    \label{fig:train_dqn_example}
\end{figure}

\begin{exercise}
    Abra el script \texttt{lunar\_lander\_test.py} (práctica de Q-learning 2, parte B). Ejecute varios tests visualizando el lander. Observe los resultados. ¿Cómo podemos evaluar correctamente el resultado del entrenamiento? ¿Cómo podemos mejorar los resultados?
\end{exercise}

Para mejorar los resultados, tendremos que ejecutar un entrenamiento con un número de episodios mayor. Además, será vital que conozcamos qué pasa cuando variamos el resto de parámetros del algoritmo.

\subsection{Análisis del código python}

\begin{exercise}
    Busque esta línea en el método \texttt{train\_networks()} de la clase     \texttt{QLearningDQN}:
    
    \texttt{self.q\_net.partial\_fit(states, target\_y)}
    
    Ejecute el código hasta esa línea en modo depuración. Anote las dimensiones de \texttt{states} y \texttt{target\_y}. Justifique estas dimensiones. ¿Por qué? ¿Qué es \texttt{states} y qué es \texttt{target\_y}?
\end{exercise}
Efectivamente, nos damos cuenta de que:
\begin{itemize}
    \item El buffer para guardar muestras alberga, como máximo 50000 estados (en la clase \texttt{QLearningDQN}).
    \begin{verbatim}
    self.replay_buffer = ReplayBuffer(capacity=50000)
    \end{verbatim}
    \item De este buffer, se muestrean aleatoriamente 64 estados y acciones.
    \item Por tanto, la entrada $X=s_t$ tiene dimensiones $64\times8$ (pues 8 es la dimensión del estado de entrada).
    \item La salida tiene dimensiones $64\times4$, pues hay un valor $y$ que se desea aproximar mediante la red neuronal ($y=r_t+\max_a{\hat{Q}(s_{t+1}, a)}$) para cada posible acción.
\end{itemize}

\begin{exercise}
    Busque estas línea en el método \texttt{test()} de la clase \texttt{QLearningDQN}:
    
    \begin{verbatim}
        while True:
            q_values = self.q_net.predict(state.reshape(1, -1))[0]
            action = np.argmax(q_values)
    \end{verbatim}
    Explique las dimensiones de \texttt{q\_values}. Explique la línea: \texttt{action = np.argmax(q\_values)}.
\end{exercise}
\textbf{Es fácil}: dado un estado $s_t$ como entrada, la red predice cuatro valores $Q(s_t, a_i)$ para cada acción $a_i$. Igual que en el método Q-learning tabular, se elige la acción correspondiente al mayor valor de $Q(s_t, a_i)$. 

\begin{exercise}
    Explique por qué se hace, en la línea de abajo 
    \texttt{target\_y = self.q\_net.predict(states)}. Nótese que en esta línea se están calculando los \textit{targets} usando la red online $Q$, cuando se ha indicado anteriormente que debería usarse la red target $\hat{Q}$ para hacer esto.
    \begin{small}    
        \begin{verbatim}
    def train_networks(self, total_steps):
        if len(self.replay_buffer) >= self.batch_size:
            states, actions, rewards, next_states, dones = self.replay_buffer.sample
                             (self.batch_size)
            target_y = self.q_net.predict(states)
            next_q_values = self.target_net.predict(next_states)
            max_next_q = np.max(next_q_values, axis=1)
            for i in range(self.batch_size):
                a = actions[i]
                r = rewards[i]
                d = dones[i]
                if d:
                    target_y[i, a] = r
                else:
                    target_y[i, a] = r + self.gamma * max_next_q[i]
            self.q_net.partial_fit(states, target_y)
        # Sync Target Network periodically
        if total_steps % self.target_update_freq == 0:
            self.target_net = copy.deepcopy(self.q_net)
        \end{verbatim}
    \end{small}
\end{exercise}

Nótese que realmente se está usando la red $\hat{Q}$ (\texttt{target\_net}) \textbf{para las acciones que se han ejecutado en cada estado}. Los valores importantes se guardan en la variable \texttt{next\_q\_values} y \texttt{max\_next\_q}. Usando estas variables se calcula la salida $y$ en las siguientes líneas:
\begin{verbatim}
    if d:
        target_y[i, a] = r
    else:
        target_y[i, a] = r + self.gamma * max_next_q[i]
\end{verbatim}

Nótese que se está calculando la salida $y$ para la columna \texttt{a} (correspondiente a la acción $a$).
Sin embargo, ¿qué valor le damos a la salida $y$ \textbf{para las acciones que no se han ejecutado en ese estado}? En este caso, se usa la red \textit{online} $Q$ para solucionar este problema. De esta manera, en las acciones que no se han ejecutado el error es cero y no se actualizarán los pesos de la red. En resumen: \texttt{target\_y = self.q\_net.predict(states)} rellena todas las salidas con la información previa de la red $Q$ y \texttt{next\_q\_values = self.target\_net.predict(next\_states)} actualiza los valores de $y$ para mejorar la estimación de la red $Q$.

%is not actually computing the Bellman target; it is a workaround for scikit-learn's API constraints.The Problem with Scikit-Learn's MLPRegressorIn deep learning frameworks like PyTorch or TensorFlow, you can compute loss using a mask so that gradients only flow back through the single action $a$ that was executed in the transition $(s, a, r, s', d)$.However, scikit-learn's partial_fit(X, Y) requires a target matrix Y of size (batch_size, 4) containing explicit target values for all 4 output neurons simultaneously.Because a single transition only tells us what happened for the one action $a$ that was taken, we have no immediate target for the other 3 actions that were not taken.How the Workaround Operates Step-by-StepTo train the network on action $a$ without corrupting or distorting the predictions for the other 3 actions, the code does the following:Step 1: Pre-fill the target matrix with current predictionsPythontarget_y = self.q_net.predict(states)
%This populates all 4 columns of target_y with the online network's current predictions for every action:$$target\_y[i] = [Q(s, a_0), Q(s, a_1), Q(s, a_2), Q(s, a_3)]$$Step 2: Overwrite ONLY the chosen action using the Target NetworkPythonnext_q_values = self.target_net.predict(next_states)
%max_next_q = np.max(next_q_values, axis=1)

%# Overwrite ONLY column 'a'
%target_y[i, a] = r + self.gamma * max_next_q[i]
%Here, the value for action $a$ is replaced with the true Bellman Target computed using target_net.Step 3: Train the network via partial_fitWhen scikit-learn evaluates the Mean Squared Error loss $(y_{\text{target}} - y_{\text{pred}})^2$ across all 4 output neurons:For the action taken ($a$):$$\text{Loss}_a = \left( \left[ r + \gamma \max_{a'} Q_{\text{target}}(s', a') \right] - Q_{\text{online}}(s, a) \right)^2 \neq 0$$This produces a non-zero gradient, updating the weights for action $a$.For unchosen actions ($a_{\text{unchosen}}$):$$\text{Loss}_{\text{unchosen}} = \left( Q_{\text{online}}(s, a_{\text{unchosen}}) - Q_{\text{online}}(s, a_{\text{unchosen}}) \right)^2 = 0$$This produces a zero gradient, ensuring the network's knowledge about unchosen actions remains untouched.self.q_net.predict(states) is used strictly to set up a "zero-loss mask" for the unchosen actions, ensuring partial_fit only modifies the Q-value of the action that was actually performed.


%\begin{exercise}
%\textbf{Exploración y explotación}. Nos interesa entender qué ocurre durante el aprendizaje cuando variamos $\epsilon$ de formas diferentes. Usando estos parámetros:
%\begin{itemize}
%    \item $\gamma=0.99$
%    \item MLP de $64\times64$.
%    \item 500 episodios de entrenamiento.
%\end{itemize}
%Probaremos a cambiar $\epsilon$ de diferentes maneras:
%\begin{itemize}
%    \item Un cambio rápido de $\epsilon$ de $1.0$ to $0.01$ durante el 20\% del entrenamiento.
%    \item Un cambio lento de $\epsilon$ de $1.0$ to $0.01$ durante el 95\% del entrenamiento.
%    \item Un valor fijo de $\epsilon=0.5$ durante todo el entrenamiento.
%\end{itemize}

%Revise y compare los resultados para los tres casos y conteste las preguntas siguientes:
%\begin{itemize}
%    \item ¿Una bajada muy rápida de $\epsilon$ (de 1 a 0.01 en el 20\% de los episodios) permite al agente aterrizar correctamente?
%    \item ¿Una bajada lenta de $\epsilon$ durante el 95\% de los episodios produce un aprendizaje correcto?
%    \item ¿Mantener un valor alto de $\epsilon$ durante todo el entrenamiento permite al agente aprender una política de aterrizaje correcta?
%\end{itemize}
%\end{exercise}

%Your result feels counter-intuitive because standard reinforcement learning intuition suggests that more exploration leads to a better policy. However, LunarLander's mechanics heavily reward fine-grained stabilization and precise control over broad spatial exploration.Dense, Continuous FeedbackUnlike sparse-reward environments (like a maze where the agent only gets a reward upon reaching the exit), LunarLander provides immediate, continuous rewards based on position, angle, velocity, and leg contact. The agent doesn't need hundreds of episodes of random actions to "find" the goal—it receives clear vector signals from step one, allowing it to map basic physics very quickly.Precision Over RandomnessLanding requires delicate micro-adjustments. When $\epsilon$ is high, random action selection introduces severe noise. Firing a main or side thruster randomly at a critical moment near the ground will instantly flip or crash the lander. Fast decay drops $\epsilon$ to $0.01$ early, preventing random exploratory actions from ruining otherwise perfect descent trajectories.High Episode Count in Steps200 episodes in LunarLander equal roughly $40,000$ to $80,000$ environment timesteps. That is more than enough transition data to sufficiently fill a replay buffer with exploratory physics samples. Decaying early does not mean under-exploring; it simply means transition to refinement happens at an appropriate volume of experience.Replay Buffer Trajectory QualityDecaying to $0.01$ by episode 200 gives your Q-network 800 episodes (80% of total training) to train on high-quality, near-deterministic flight trajectories. Slower decays (like 950 episodes) keep the replay buffer constantly polluted with erratic, noisy transitions, slowing down optimal policy convergence.Why Constant Epsilon FailsConstant values demonstrate why persistent noise breaks control tasks:$\epsilon = 0.95$ / $0.50$: The agent acts randomly most of the time, making controlled descent impossible.$\epsilon = 0.20$: Taking a completely random action 20% of the time means an error occurs every 5 timesteps. In a delicate physics landing, a 20% mistake rate guarantees high crash rates, artificially capping maximum performance.

\subsection{Cambios en los parámetros básicos del algoritmo}
A continuación, ``jugaremos'' con el algoritmo y nos fijaremos en los resultados que se obtienen al cambiar los hiperparámetros. \textbf{Nota}: en el ámbito de la Inteligencia Artificial es muy importante mantener un registro de los resultados que se obtienen para cada conjunto de parámetros. Así pues, para facilitar el trabajo al estudiante se proporciona un método \texttt{save\_results} para guardar los resultados de train o de test en un fichero \texttt{.json} (en el directorio \texttt{results}). Posteriormente, se pueden ver los resultados de experimentos diferentes en formato html si ejecutamos el script \texttt{plot\_results.py}.

%\begin{exercise}
%Abra el script \texttt{lunar\_lander\_train\_multi.py}. Verá una función \texttt{set\_experiments} donde se pueden escribir, como diccionarios, los parámetros de un conjunto de simulaciones. De esta manera se pueden programar simulaciones y visualizar los resultados más tarde. Al finalizar cada simulación, se guarda un fichero \texttt{.json} con los resultados de train y de test.
%\end{exercise}


\begin{exercise}
    \textbf{Se desea ver qué ocurre al cambiar $\gamma$ (\textit{discount factor})}, con $\gamma=0.1, 0.3, 0.5, 1.0$, denominándose:
    \begin{itemize}
        \item $\gamma = 0.1$: miope o cortoplacista.
        \item $\gamma = 0.99$: visionario o largoplacista.
        \item $\gamma = 1.0$: visionario o largoplacista sin factor de descuento.
    \end{itemize}
    Utilice el script \texttt{lunar\_lander\_train\_multi.py}. Modifique la función \texttt{set\_experiments} para utilizar los parámetros deseados. Use 1000 pasos de entrenamiento y observe la recompensa media y la varianza. Tras cada experimento, se guarda un fichero \texttt{json}. Ejecute \texttt{plot\_results.py} y observe qué resultados de entrenamiento y test se obtienen (se abren dos pestañas en el navegador web).  Responda:
    \begin{itemize}
        \item ¿Un valor bajo de $\gamma$ permite al \textit{lander} orientarse correctamente y reducir el consumo de combustible?
        \item ¿Qué ocurre con un valor de $\gamma$ cercano a 0.99?
        \item ¿Qué ocurre con un valor de $\gamma=1.0$?
    \end{itemize}
     %Observe the physical behavior of both agents during early-to-mid training steps.Questions to Answer:How does a low $\gamma$ alter the agent's strategy regarding fuel consumption and early orientation maneuvers?Mathematically, what happens to the return calculation $G_t$ when $\gamma \rightarrow 0$ vs. $\gamma \rightarrow 1$?
\end{exercise}
Nos damos cuenta de que la implementación con $\gamma=1.0$ no converge correctamente. En concreto, empieza a divergir cuando $\epsilon$ disminuye. La razón tiene que ver con la estimación de la red neuronal Q. Considere el caso en el que $\gamma = 1.0$. La recompensa acumulada $Q(s,a)=r_t+\gamma r_{t+1}+\gamma^2 r_{t+2} ...=Q(s,a)=r_t+ r_{t+1}+ r_{t+2} ...$. De esta manera los valores de recompensa total crecen sin límite en episodios de aproximadamente 200 pasos. Esto da una gran varianza en los valores que debe aproximar la red neuronal y típicamente no converge. Se puede obtener un efecto parecido si fijamos $\gamma=0.99$ con la que la red sí que converge.

Al tener diferentes resultados cambiando $\gamma$ nos damos cuenta de que un \textit{discount factor} bajo permite al módulo aprender a orientarse correctamente, acercarse a la posición $(0,0)$ y mantener una velocidad baja, pero no le permite aprender a aterrizar correctamente y recibir la recompensa final. Con $\gamma=1.0$ el agente es capaz de aprender correctamente a aterrizar en el lugar designado.

Ciertamente, la realización de los experimentos consume bastante tiempo. Para que los resultados sean concluyentes, se necesita un mayor número de episodios de entrenamiento. Además, es importante darse cuenta de que el sistema es, en extremo aleatorio, con lo que el resultado del aprendizaje depende fuertemente de los estados iniciales que haya tomado el \textit{lander} durante el primer paso de los episodios de entrenamiento y de las acciones aleatorias que se han llevado a cabo. En consecuencia, \textbf{es necesario calcular la media entre los resultados de entrenamientos diferentes}.





\begin{exercise}
De los experimentos anteriores, ya debería contar con un valor de $\gamma$ adecuado. Con el valor óptimo de este parámetro nos planteamos variar el tamaño de la red MLP que aproxima la tabla Q. Compare los resultados de estas tres situaciones:
\begin{itemize}
    \item Infradimensionada: \texttt{hidden\_layer\_sizes=(8, 8)}.
\item Estándar: \texttt{hidden\_layer\_sizes=(64, 64)}.
\item Sobredimensionada: \texttt{hidden\_layer\_sizes=(256, 256, 256)}.
\end{itemize}

Utilice el script Utilice el script \texttt{lunar\_lander\_train\_multi.py}. Modifique la función \texttt{set\_experiments} para utilizar los parámetros deseados. Observe qué resultados de entrenamiento y test se obtienen (plot\_results.py). Use 1000 pasos de entrenamiento y observe la recompensa media y la varianza. ¿Cuál de las redes MLP permite obtener mejores resultados?
\end{exercise}

%Measure two metrics for each setup: total time per 100 episodes (CPU execution speed) and the average episode score after 300 episodes.

%Questions to Answer:

%Why is the tiny network incapable of learning complex flight physics?

%Does the largest network converge faster in terms of episode count? Why might a larger network slow down overall wall-clock training time when using CPU-bound models like scikit-learn?
\subsection{Cambios en los parámetros avanzados del algoritmo}

Hay unos parámetros de la clase \texttt{QLearningDQN} que no resultan obvios a primera vista:
\begin{itemize}
    \item \texttt{replay\_buffer\_size}: es el tamaño del buffer que almacena las muestras de estados, acciones y recompensas que han ocurrido durante el entrenamiento:
\begin{verbatim}
self.replay_buffer = ReplayBuffer(capacity=params.get('replay_buffer_size', 50000))
\end{verbatim}
    Si se disminuye el tamaño de este buffer, estaremos usando solamente la historia reciente del estado y acciones (es un buffer FIFO). El algoritmo requiere que haya una buena variedad de estados para poder muestrear estados que estén alejados en el tiempo.
    
    \item \texttt{target\_update\_freq}: es la frecuencia en pasos con la que se sobreescribe la red neuronal target con los pesos de la red online $\hat{Q}=Q$. Por defecto \texttt{target\_update\_freq=250}, nos disponemos a ver qué ocurre cuando bajamos este valor.
    \item \texttt{batch\_size}: es el número de datos $X=(s_t)$ que se utilizan en cada iteración SGD de la red $Q$. Un número bajo puede conllevar problemas de convergencia.
\end{itemize}

\begin{exercise}
El uso de un buffer intermedio para estimar $Q$ (red online) se justifica para evitar que la regresión esté sesgada hacia los últimos valores del estado. Utilice el script \texttt{lunar\_lander\_train\_multi.py}. Modifique la función \texttt{set\_experiments} para variar el parámetro \texttt{replay\_buffer\_size=1, 100, 1000, 10000, 50000}. Observe los resultados usando \texttt{plot\_results.py}.
\end{exercise}

\begin{exercise}
El uso de dos redes diferenciadas $Q$ (online) y $\hat{Q}$ (target) se justifica para estabilizar la estimación de $Q$. En este ejercicio se busca demostrar que, en efecto, es necesario mantener una copia estática ($\hat{Q}$, target) durante un número de pasos determinado. Utilice el script \texttt{lunar\_lander\_train\_multi.py}. Modifique la función \texttt{set\_experiments} para variar el parámetro \texttt{target\_update\_freq=1, 10, 100, 250 500, 1000}. Observe los resultados usando \texttt{plot\_results.py}.
\end{exercise}

\begin{exercise}
En cada paso del algoritmo, se realiza una única iteración por SGD de la red neuronal. El parámetro \texttt{batch\_size} es el número de estados de entrada $X=s_t$ para los que se estimará una salida mediante la red. Estos \texttt{batch\_size} elementos se muestrean del buffer intermedio ReplayBuffer. Utilice el script \texttt{lunar\_lander\_train\_multi.py}. Modifique la función \texttt{set\_experiments} para variar el parámetro \texttt{batch\_size=1, 10, 64, 200 500, 1000}. Observe los resultados usando \texttt{plot\_results.py}.
\end{exercise}


\section{Proyectos e ideas}
%\begin{exercise}
%    ¿Cómo adaptarías este algoritmo DQN para resolver el problema del bipedal walker?
% Nada, pregunta compleja, pues necesitas dos redes neuronales
%\end{exercise}

\begin{exercise}
    \textcolor{red}{Juanjo: me estoy peleando por ver si en DQN puedo encontrar algo un pelín más eficiente para que se tarde menos en hacer las pruebas. De todos los ejercicios que se plantean, realmente no creo que dé tiempo a mucho. Estoy pensando realmente en dejar los ficheros json en un directorio y que los copien para visualizarlos. Además, los resultados son muy variables, pues es todo altamente aleatorio: hay que repetir el experimento con gamma=0.8 unas 10 veces para que tenga sentido el resultado.}
\end{exercise}

\begin{exercise}
    \textcolor{red}{Juanjo: se me ocurre que se podría montar un concurso. Cada grupo de 2 estudiantes envía su versión final de DQN del lunar lander. Se simulan 1000 episodios y gana el que tenga el mejor average reward.}
\end{exercise}

%In contrast to TD-Gammon and similar online approaches, we utilize a technique known as expe-
%rience replay [13] where we store the agent’s experiences at each time-step, et = (st , at , rt , st+1 )
%in a data-set D = e1 , ..., eN , pooled over many episodes into a replay memory. During the inner
%loop of the algorithm, we apply Q-learning updates, or minibatch updates, to samples of experience,
%e ∼ D, drawn at random from the pool of stored samples. After performing experience replay,
%the agent selects and executes an action according to an -greedy policy. Since using histories of
%arbitrary length as inputs to a neural network can be difficult, our Q-function instead works on fixed
%length representation of histories produced by a function φ. The full algorithm, which we call deep
%Q-learning, is presented in Algorithm 1.

%This approach has several advantages over standard online Q-learning [23]. First, each step of
%experience is potentially used in many weight updates, which allows for greater data efficiency.

\bibliographystyle{alpha}
\bibliography{sample}

\end{document}
